\section{Introducci\'on}
\label{sec:introduccion}

La transferencia de calor por conducci\'on interviene en procesos de calentamiento, enfriamiento y disipaci\'on t\'ermica. La conductividad caracteriza la capacidad de los materiales para transportar calor y condiciona el dise\~no de sistemas de gesti\'on t\'ermica \citep{chung2001materials,zheng2021advances}. En dominios bidimensionales, la distribuci\'on de temperatura depende de la geometr\'ia, las propiedades del material, las fuentes internas y las condiciones de frontera. Cuando el dominio no permite una soluci\'on anal\'itica directa, una discretizaci\'on espacial conduce a un problema algebraico que puede resolverse num\'ericamente.

En sistemas electr\'onicos, parte de la energ\'ia consumida por los componentes se transforma en calor. Los modelos t\'ermicos de procesadores han representado el flujo entre bloques funcionales y el encapsulado mediante redes compactas; los estudios de placas impresas han incorporado fuentes cercanas, convecci\'on, radiaci\'on y estructuras multicapa \citep{skadron2003hotspot,harvest2007devices,moniervinard2015pcb,zhang2024pcb}. Estos trabajos describen fen\'omenos que no se incluyen en el presente modelo: la CPU y la GPU se idealizan como regiones rectangulares con generaci\'on uniforme, dentro de una placa homog\'enea bidimensional y con temperatura de contorno fija. Por tanto, las temperaturas obtenidas caracterizan el sistema matem\'atico definido y no mediciones del funcionamiento de una laptop.

Los problemas bidimensionales de conducci\'on se han resuelto mediante diferencias finitas, formulaciones de volumen finito y m\'etodos adaptativos para ecuaciones de Poisson. En una placa cuadrada estacionaria, \citet{santos2022finite} compararon el M\'etodo de Diferencias Finitas (MDF) con una soluci\'on anal\'itica y reportaron errores relativos menores que 0.003 sin generaci\'on y de hasta 0.03 con generaci\'on; ese resultado respalda el estudio num\'erico de la conducci\'on, aunque no verifica la placa ni las fuentes de este proyecto. \citet{oladejo2007tool} combinaron el MDF con Gauss--Seidel para una placa cuadrada, mientras \citet{settle2013compact} analizaron la precisi\'on de esquemas compactos de cinco puntos para Poisson. Otros trabajos desarrollaron discretizaciones para dominios adaptativos y formulaciones de volumen finito \citep{chen2007supra,raeli2018hierarchical,batty2017cell}. Estos antecedentes sit\'uan el stencil de cinco puntos usado aqu\'i dentro de una familia de m\'etodos para ecuaciones el\'ipticas, pero no sustituyen su verificaci\'on para la configuraci\'on estudiada.

La selecci\'on del solucionador depende de la estructura del sistema discreto. Gradiente Conjugado se formula para matrices sim\'etricas definidas positivas; \citet{zhang2001} lo aplicaron en una formulaci\'on de transporte t\'ermico de microescala, distinta del modelo estacionario aqu\'i considerado. Los an\'alisis de m\'etodos iterativos para sistemas dispersos derivados de problemas de Poisson ofrecen contexto para comparar la convergencia, sin anticipar el desempe\~no en esta malla particular \citep{xu2011poisson}. Como alternativas espaciales, \citet{akeremale2020finite} estudiaron una placa rectangular mediante elementos finitos y \citet{reutskiy2016meshless} aplicaron funciones de base radial a la conducci\'on estacionaria en medios anis\'otropos y no homog\'eneos. La verificaci\'on de una soluci\'on de conducci\'on requiere distinguir entre concordancia algebraica y exactitud espacial; \citet{beck2004verification} presentaron soluciones de referencia para s\'olidos rectangulares con calentamiento parcial, bajo condiciones que no coinciden por completo con las del proyecto.

Aunque estos estudios muestran m\'etodos aplicables a problemas t\'ermicos, no eval\'uan conjuntamente el campo de una placa idealizada con dos fuentes rectangulares, la convergencia algebraica de Gauss--Seidel y Gradiente Conjugado, el refinamiento espacial y la sensibilidad a las potencias. El estudio se centra en esa comparaci\'on reproducible. Se emplea el MDF para discretizar la ecuaci\'on estacionaria y se comparan ambos solucionadores sobre el mismo sistema disperso, separando la respuesta del modelo de los errores algebraicos y de discretizaci\'on.

El objetivo general de la investigaci\'on es modelar num\'ericamente la distribuci\'on de temperatura en una placa electr\'onica idealizada con fuentes t\'ermicas localizadas asociadas con la CPU y la GPU, con el prop\'osito de analizar su comportamiento t\'ermico y comparar el desempe\~no de los m\'etodos de Gauss--Seidel y Gradiente Conjugado en la resoluci\'on del sistema discretizado mediante diferencias finitas. De manera espec\'ifica, se busca: (i) formular e implementar el modelo num\'erico bidimensional de conducci\'on estacionaria mediante diferencias finitas, considerando la ecuaci\'on gobernante, las condiciones de frontera y las fuentes t\'ermicas localizadas asociadas con la CPU y la GPU; (ii) evaluar el comportamiento del modelo bajo un conjunto base de par\'ametros mediante la distribuci\'on de temperatura, la temperatura m\'axima y las m\'etricas de convergencia de los m\'etodos num\'ericos; (iii) analizar la robustez del modelo mediante la variaci\'on de la generaci\'on de calor asociada con la CPU y la GPU y de los par\'ametros de discretizaci\'on espacial; y (iv) aplicar el modelo formulado para comparar el desempe\~no de Gauss--Seidel y Gradiente Conjugado mediante el error o residuo num\'erico, el n\'umero de iteraciones y el costo computacional requerido para alcanzar un criterio de convergencia establecido. El art\'iculo presenta el modelo y los procedimientos en Metodolog\'ia, los hallazgos en Resultados, su interpretaci\'on cr\'itica en Discusi\'on y la s\'intesis final en Conclusiones.
